\documentclass[11pt,a4paper]{article}

% --- Packages ---
\usepackage[utf8]{inputenc}
\usepackage[english]{babel}
\usepackage{amsmath, amssymb, amsfonts, bm}
\usepackage{geometry}
\usepackage{graphicx}
\usepackage{subcaption}
\usepackage{booktabs}
\usepackage{hyperref}
\usepackage{cite}
\usepackage{float}
\usepackage{xcolor}

\geometry{a4paper, top=2.5cm, bottom=2.5cm, left=2.5cm, right=2.5cm}

\hypersetup{
    colorlinks=true,
    linkcolor=blue!70!black,
    citecolor=green!50!black,
    urlcolor=cyan!70!black
}

% --- Macros ---
\newcommand{\bq}{\bm{q}}
\newcommand{\bdq}{\dot{\bm{q}}}
\newcommand{\bddq}{\ddot{\bm{q}}}
\newcommand{\btau}{\bm{\tau}}
\newcommand{\bM}{\bm{M}}
\newcommand{\bC}{\bm{C}}
\newcommand{\bG}{\bm{G}}
\newcommand{\bB}{\bm{B}}
\newcommand{\Real}{\mathbb{R}}

% ============================================================
\begin{document}

\begin{titlepage}
    \centering
    \vspace*{2cm}
    {\Large\bfseries Field and Service Robotics --- Final Technical Report\par}
    \vspace{1.5cm}
    {\LARGE\bfseries Modeling and Multi-Strategy Control of a\\[6pt]
    Variable-Length Wheeled Inverted Pendulum Robot\par}
    \vspace{1cm}
    {\large RoTino: a Wheeled Biped Balancing Platform\par}
    \vspace{2cm}
    {\large Aldo Rosicarelli\par}
    \vspace{0.5cm}
    {\large Academic Year 2025--2026\par}
    \vfill
    {\large Field and Service Robotics\par}
\end{titlepage}

% ============================================================
\tableofcontents
\newpage

% ============================================================
\section{Introduction}
\label{sec:intro}

Wheeled biped robots occupy a particularly interesting niche in the landscape of field and service
robotics: they combine the high mobility and energy efficiency of wheeled platforms with the
postural versatility of legged systems, allowing them to adjust their centre of mass (CoM) height
to negotiate uneven terrain, step over obstacles, or adopt a more compact stance.  The archetype
is the family of Variable-Length Wheeled Inverted Pendulum (VL-WIP) robots --- two-wheeled
self-balancing machines whose legs can change length to modulate the height at which the
upper-body mass is carried.

The project selected for this course is \textbf{RoTino}, a 2D simulation of a wheeled biped
modelled after the platform described by Cui et al.~\cite{cui2022}.  The choice is motivated by
several factors:

\begin{itemize}
    \item \textbf{Rich nonlinear underactuated dynamics.}  A four-DOF planar chain with only
          three independent actuators (wheel torque, knee, hip) requires careful treatment of
          the actuation matrix and its singularities.
    \item \textbf{Cascade control hierarchy.}  The natural time-scale separation between balance
          (fast) and locomotion (slow) makes the system an ideal testbed for hierarchical control
          strategies of increasing sophistication.
    \item \textbf{Comparison across control paradigms.}  Three qualitatively different
          controllers --- a classical cascade PID, an LQR coupled with a Model Predictive
          Controller (MPC), and a nonlinear Sliding Mode Controller (SMC) --- are designed,
          implemented, and benchmarked on the same model, enabling a principled comparison of
          performance, robustness, and implementation complexity.
    \item \textbf{Software-in-the-loop validation.}  The MATLAB/Simulink environment is used for
          controller prototyping, and a ROS~2 + Gazebo pipeline is used for closed-loop
          simulation of the complete robot, bridging the gap between symbolic derivation and
          practical deployment.
\end{itemize}

The report is structured as follows.  Section~\ref{sec:topology} analyses the topology and
degrees of freedom.  Section~\ref{sec:model} derives the Lagrangian dynamic model.
Sections~\ref{sec:pid}, \ref{sec:lqr_mpc}, and \ref{sec:smc} describe the three control
strategies.  Section~\ref{sec:results} presents the simulation results.
Section~\ref{sec:conclusions} concludes.

% ============================================================
\section{Topology Analysis and Degrees of Freedom}
\label{sec:topology}

\subsection{Mechanical Architecture}

The robot (Figure~\ref{fig:architecture}) is a planar kinematic chain comprising four rigid
bodies:

\begin{enumerate}
    \item \textbf{Wheel} ($r = 0.127$~m, $m_w = 3.5$~kg, $I_w = 0.1$~kg\,m$^2$): rolling
          without slipping on the ground; the rolling constraint eliminates one translational DOF.
    \item \textbf{Shank link} ($l_1 = 0.45$~m, $m_1 = 1.2$~kg): connects the wheel axle to
          the knee joint.
    \item \textbf{Thigh link} ($l_2 = 0.45$~m, $m_2 = 5.3$~kg): connects the knee to the hip
          joint.
    \item \textbf{Torso} ($l_3 = 0.35$~m, $m_3 = 60.0$~kg, wheel track $d = 0.63$~m): the
          upper body, connected to the hip; its orientation evolves freely and is
          \emph{not} independently actuated.
\end{enumerate}

\begin{figure}[H]
    \centering
    \includegraphics[width=0.38\textwidth]{figures/mechanical_architecture.png}
    \caption{Schematic of the RoTino platform: torso, hip joint, thigh link, knee joint, shank
    link, wheel and wheel joint.}
    \label{fig:architecture}
\end{figure}

\subsection{Constraint and DOF Count}

A planar system of $n_b = 4$ bodies has $3 n_b = 12$ unconstrained DOFs.  Table~\ref{tab:dof}
summarises the constraints that reduce this count to 4.

\begin{table}[H]
    \centering
    \caption{Constraint analysis for the planar RoTino chain.}
    \label{tab:dof}
    \begin{tabular}{llcc}
        \toprule
        Constraint & Description & Equations & Holonomic? \\
        \midrule
        Wheel--ground contact & Fixed contact height $y = r$ & 1 & Yes \\
        Pure rolling & $\dot{x} = r\,\dot{\phi}_w$ & 1 & No (Pfaffian) \\
        Wheel--shank revolute (ankle) & 2 translational DOFs locked & 2 & Yes \\
        Shank--thigh revolute (knee) & 2 translational DOFs locked & 2 & Yes \\
        Thigh--torso revolute (hip) & 2 translational DOFs locked & 2 & Yes \\
        \midrule
        \textbf{Total constraints} & & \textbf{8} & \\
        \textbf{Resulting DOF} & & \textbf{4} & \\
        \bottomrule
    \end{tabular}
\end{table}

The generalised coordinates, chosen to match Cui et al.\ (Table~1), are:
\begin{equation}
    \bq = \begin{bmatrix} q_w & q_1 & q_2 & q_3 \end{bmatrix}^T
\end{equation}
where $q_w$ is the wheel rotation (encoding horizontal displacement $s = R_w q_w$ via the rolling
constraint), $q_1$ is the absolute shank angle from the vertical, $q_2$ is the absolute thigh
angle, and $q_3$ is the absolute torso angle.

\subsection{Actuation, Underactuation, and Ankle Constraint}

Three independent actuators are available:
\begin{itemize}
    \item $\tau_w$: wheel torque (drives $q_w$, i.e.\ forward displacement);
    \item $\tau_1$: knee torque (shank-to-thigh);
    \item $\tau_2$: hip torque (thigh-to-torso).
\end{itemize}

The torso angle $q_3$ is \textbf{unactuated}: it evolves from the coupled dynamics, making the
system underactuated with one degree of underactuation.  Controlling $q_3$ indirectly --- via
the centroidal coupling to the wheel acceleration --- is the fundamental challenge of the balance
problem.

The ankle angle satisfies the algebraic constraint $q_0 = q_1 - q_2 + q_3$, so it is not
an independent state but is monitored as a diagnostic variable.

% ============================================================
\section{Dynamic Model}
\label{sec:model}

\subsection{Lagrangian Derivation}

The absolute barycentre positions (wheel--ground contact as origin) are:
\begin{align}
    x_1 &= s + \tfrac{l_1}{2}\sin\phi_s, \quad
    y_1 = R_w + \tfrac{l_1}{2}\cos\phi_s \\[3pt]
    x_2 &= s + l_1\sin\phi_s + \tfrac{l_2}{2}\sin\phi_t, \quad
    y_2 = R_w + l_1\cos\phi_s + \tfrac{l_2}{2}\cos\phi_t \\[3pt]
    x_3 &= s + l_1\sin\phi_s + l_2\sin\phi_t + \tfrac{l_3}{2}\sin q_3, \quad
    y_3 = R_w + l_1\cos\phi_s + l_2\cos\phi_t + \tfrac{l_3}{2}\cos q_3
\end{align}
where $\phi_s = q_1 - q_2 + q_3$ and $\phi_t = q_3 - q_2$.

The kinetic and potential energies are:
\begin{align}
    T &= \tfrac{1}{2}(m_w + I_w/R_w^2)\dot{s}^2 + \sum_{k=1}^{3}\!\left[\tfrac{1}{2}m_k\|\dot{\bm{r}}_k\|^2 + \tfrac{1}{2}I_k\dot{\phi}_k^2\right] \\
    V &= m_w g R_w + m_1 g y_1 + m_2 g y_2 + m_3 g y_3
\end{align}

The Euler--Lagrange equations $\frac{d}{dt}\frac{\partial L}{\partial \bdq} -
\frac{\partial L}{\partial \bq} = \bB\btau$ give the standard structured form:
\begin{equation}
    \bM(\bq)\bddq + \bC(\bq,\bdq)\bdq + \bG(\bq) = \bB\btau
    \label{eq:eom}
\end{equation}
with $\bM \in \Real^{4\times4}$ (symmetric, positive definite), $\btau = [\tau_w; \tau_1;
\tau_2]$, and:
\begin{equation}
    \bB = \begin{bmatrix}
        -1 & 1 & 0 \\
         1 & 0 & 0 \\
         1 & 0 & 1 \\
        -1 & 0 & 0
    \end{bmatrix} \in \Real^{4\times3}
\end{equation}
The matrices $\bM$, $\bC\bdq$, $\bG$ are derived symbolically in MATLAB via the Symbolic Toolbox
and exported as optimised numerical functions for real-time evaluation in Simulink and in the
ROS~2 Python nodes.

\subsection{Equilibrium Verification}

At the upright equilibrium ($\bq = \bm{0}$, $\bdq = \bm{0}$) the gravitational vector vanishes:
$\bG|_{\bq=\bm{0}} = \bm{0}$, which is verified numerically.

\subsection{Decoupled VL-WIP Model for LQR/MPC}

Treating the upper body as a lumped mass at a variable height $l$ above the wheel axle and
linearising around the upright, one obtains the VL-WIP state-space model for LQR design:
\begin{equation}
    \dot{\bm{x}} = A\bm{x} + B_u\bm{u}, \qquad
    \bm{x} = [s,\ \theta,\ \phi_w,\ \dot{s},\ \dot{\theta},\ \dot{\phi}_w]^T
\end{equation}

For the MPC the upper body is further decoupled.  The horizontal subsystem (states $[s, \dot{s},
\Delta s]$, $\Delta s$ the CoM offset) and the vertical subsystem (states $[z, \dot{z}]$) are
governed by:
\begin{equation}
    A_{MPC} = \begin{bmatrix}
        0 & 1 & 0 & 0 & 0 \\
        0 & 0 & 0 & 0 & g/h_0 \\
        0 & 0 & 0 & 1 & 0 \\
        0 & 0 & 0 & 0 & 0 \\
        0 & 0 & 0 & 0 & 0
    \end{bmatrix},\ 
    B_{MPC} = \begin{bmatrix}
        0 & 0 \\ 0 & 0 \\ 0 & 0 \\ 0 & 1/m_b \\ 1 & 0
    \end{bmatrix}
\end{equation}
with inputs $[\dot{\Delta s},\,F_z]^T$, nominal height $h_0 = 0.6$~m, and total upper-body mass
$m_b = 66.5$~kg.

\subsection{CoM Kinematics and Virtual Model Controller}

The horizontal and vertical CoM coordinates relative to the wheel axle are:
\begin{align}
    S_C &= C_a\sin\phi_s + C_b\sin\phi_t + C_c\sin q_3 \\
    Z_C &= C_a\cos\phi_s + C_b\cos\phi_t + C_c\cos q_3
\end{align}
with $C_a = (m_1 l_1/2 + (m_2+m_3)l_1)/m_b$, $C_b = (m_2 l_2/2 + m_3 l_2)/m_b$,
$C_c = m_3 l_3/(2m_b)$.  The pitch angle of the upper body and its analytic Jacobian are:
\begin{equation}
    \theta_b = \arctan2(S_C,\,Z_C), \qquad
    J_\theta = \frac{Z_C J_S - S_C J_Z}{S_C^2 + Z_C^2}
\end{equation}
These expressions are used by the SMC to compute the sliding-surface gradient analytically at
each timestep, eliminating the need for numerical differentiation.

% ============================================================
\section{Cascade PID Controller}
\label{sec:pid}

\subsection{Architecture}

The simplest controller is a two-level cascade PID.

\textbf{Outer loop (velocity controller).}  A PI controller maps the wheel velocity error
$e_v = \dot{q}_w^* - \dot{q}_w$ to a shank-angle offset $\delta q_1$:
\begin{equation}
    \delta q_1 = -(K_{P,v}\,e_v + K_{I,v}\!\int e_v\,dt),
    \qquad K_{P,v} = 0.08,\ K_{I,v} = 0.06
\end{equation}
The negative sign reflects the non-minimum phase nature of the plant: to accelerate forward the
robot must first tilt its axle backward.

\textbf{Inner loop (balance controller).}  A PD loop stabilises the shank angle to the modified
reference $q_1^* + \delta q_1$:
\begin{equation}
    \tau_w = K_{P,w}(q_1^* + \delta q_1 - q_1) + K_{D,w}(0 - \dot{q}_1),
    \quad K_{P,w} = 400,\ K_{D,w} = 80\ \text{[Nm/rad]}
\end{equation}

\textbf{Posture controller.}  Two independent PD loops maintain knee and hip at their
inverse-kinematic targets (desired CoM height $h_{des} = 0.75$~m, solved via \texttt{fsolve}):
\begin{equation}
    \tau_k = K_{P,j}(q_k^* - q_k) + K_{D,j}(0 - \dot{q}_k),
    \quad K_{P,j} = 800,\ K_{D,j} = 50\ \text{[Nm/rad]}
\end{equation}

\subsection{Design Choices and Limitations}

Gains are tuned empirically without explicit use of the dynamic model.  The principal limitation
is the absence of any compensation for the nonlinear coupling between the balance and posture
loops.  When the robot deviates significantly from the linearisation point the fixed gains can
no longer prevent divergence, as evidenced in the simulation results.

% ============================================================
\section{LQR + MPC Controller}
\label{sec:lqr_mpc}

\subsection{TV-LQR Balance Loop}

The linearised VL-WIP model is scheduled as a function of leg length $l$ (Time-Varying LQR,
TV-LQR).  The state and input cost matrices are:
\begin{equation}
    Q = \operatorname{diag}(30,\ 400,\ 80,\ 15,\ 6,\ 2)
\end{equation}
The pitch angle ($Q_{22} = 400$) is penalised most to ensure rapid tilt recovery.  The input
cost:
\begin{equation}
    R = T_{cd}^T \operatorname{diag}(R_{cm},\ R_{diff})\,T_{cd},
    \quad R_{cm} = 2.0,\quad R_{diff} = 100.0
\end{equation}
heavily penalises the differential wheel torque (yaw) relative to the common-mode torque
(balance), keeping the yaw bandwidth below the $\approx 12.5$~Hz torsional resonance of the
stance legs.

\subsection{Upper-Body MPC}

A linear MPC with horizon $N = 25$ steps at $\Delta t = 0.02$~s (100~Hz) plans the CoM offset
$\Delta s$ and vertical force $F_z$ over a 0.5~s window.  The QP is solved at each MPC tick
using a Python solver.

\subsection{VMC Leg Controller}

The MPC outputs $(\Delta s, z_{ref}, F_z)$ drive hip and knee torques through a Virtual Model
Controller.  In stance, only a horizontal spring realises $\Delta s$; the vertical DOF is driven
by the feedforward force $F_z$ from the MPC:
\begin{equation}
    K_P = \operatorname{diag}(1500,\ 0),\quad K_D = \operatorname{diag}(40,\ 15)\
    \text{[N/m, N\,s/m]}
\end{equation}

\subsection{Kalman State Estimator}

A linear Kalman filter at 500~Hz fuses IMU acceleration, wheel odometry, and leg kinematics to
produce $(s,\dot{s},\theta,\dot{\theta})$.  Process noise: $Q_{acc} = 0.5$; measurement noise:
$R_{pos} = 10^{-4}$, $R_{vel} = 10^{-3}$.

% ============================================================
\section{Sliding Mode Controller}
\label{sec:smc}

\subsection{Architecture}

The SMC is a nonlinear cascade with two time-scales:
\begin{enumerate}
    \item \textbf{Outer loop --- velocity to pitch reference.}  A PI controller with
          \emph{negative} gains (non-minimum phase correction) maps CoM velocity error to
          $\theta^*$:
          \begin{equation}
              \theta^* = \operatorname{sat}\!\left(K_P e_v + K_I\!\int e_v\,dt,\ \pm\theta_{max}\right),
              \quad K_P = -0.22,\ K_I = -0.06,\ \theta_{max} = 12^\circ
          \end{equation}
          Anti-windup freezes the integrator when saturation is active.
    \item \textbf{Inner loop --- three SMC channels on} $\bm{y} = [\theta_b;\,q_1;\,q_2]$
          at 500~Hz, with sliding surfaces:
          \begin{equation}
              s_k = \dot{e}_k + c_k e_k, \quad k = 1,2,3
          \end{equation}
\end{enumerate}

\subsection{Super-Twisting Pitch Channel}

Channel~1 uses the super-twisting second-order SMC algorithm, which is \emph{continuous} (no
chattering in the torque signal)~\cite{levant1993, moreno2012}:
\begin{align}
    \nu_1 &= -k_a\sqrt{|s_1|}\,\frac{s_1}{|s_1|+\varepsilon} + W \\[4pt]
    \dot{W} &= -k_b\,\frac{s_1}{|s_1|+\varepsilon}
\end{align}
with $c_1 = 6$~s$^{-1}$, $k_a = 3.5$, $k_b = 12.0$, $\varepsilon = 0.01$ (regularisation to
avoid chattering at the discontinuity).  The integrator state $W$ is frozen when the wheel
torque saturates, preventing windup under large disturbances.

\subsection{Posture Channels (Knee and Hip)}

Channels~2 and~3 use classical SMC with boundary layer (continuous approximation):
\begin{equation}
    \nu_k = -k_k\,\operatorname{sat}\!\left(\tfrac{s_k}{\phi_k}\right) - \eta_k s_k,
    \quad k = 2,3
\end{equation}
with $c_{2,3} = 12$~s$^{-1}$, $k_{2,3} = 60$, $\eta_{2,3} = 20$, $\phi_{2,3} = 0.05$~rad.
The linear term $\eta_k s_k$ provides exponential attraction inside the boundary layer.

\subsection{Partial Feedback Linearisation Inversion}

The three virtual inputs $\bm{\nu}$ are mapped to actual torques by inverting the decoupled
output dynamics:
\begin{equation}
    A_d = J_y\,\bM^{-1}\bB \in \Real^{3\times3}, \qquad
    b_d = -J_y\,\bM^{-1}(\bC\bdq + \bG) + h_y
\end{equation}
\begin{equation}
    \btau = A_d^{-1}\!\left(-b_d - c \odot \dot{\bm{e}} + \bm{\nu}\right)
\end{equation}
where $J_y = [J_{\theta_b};\,e_1^T;\,e_3^T]$ is the output Jacobian and $h_y$ collects
Lie-derivative correction terms.  The singularity indicator $\sigma = M_{22} + M_{42}$ is
monitored at runtime; $\sigma \to 0$ signals that the fully extended leg is approaching a
kinematic singularity where the wheel loses authority over pitch.

% ============================================================
\section{Results and Discussion}
\label{sec:results}

All controllers are benchmarked in two scenarios: (i)~\textbf{free balance} (released from rest,
$v^* = 0$, 26~s simulation) and (ii)~\textbf{push recovery} (50~N impulsive horizontal force
at $t = 4$~s, 20~s simulation).

\subsection{Free Balance Benchmark}

Figure~\ref{fig:dashboard} shows the overall dashboard for the free-balance scenario.

\begin{figure}[H]
    \centering
    \includegraphics[width=\linewidth]{figures/dashboard_overview.png}
    \caption{Benchmark overview: pitch angle, wheel torque command, linear velocity, and contact
    normal force for PID (red), MPC/LQR (blue), and SMC (green).}
    \label{fig:dashboard}
\end{figure}

The PID controller (red) fails to stabilise the robot.  The pitch grows into a sustained
oscillation of amplitude $\pm 15^\circ$--$18^\circ$, the wheel torque saturates at $\pm 6$~Nm, and the
contact force fluctuates between 0 and 42~N, indicating intermittent wheel lift-off.  The
velocity drifts sinusoidally at up to $\pm 3$~m/s with no tendency to settle.

In contrast, both MPC (blue) and SMC (green) stabilise the robot within $\approx 2$~s and
maintain pitch within $\pm 0.5^\circ$ thereafter, with wheel torque below 0.15~Nm --- three orders
of magnitude smaller than the PID effort.

\subsubsection{Position and Velocity Tracking}

Figure~\ref{fig:posvel} confirms that under MPC and SMC the robot stays essentially stationary
($|x| < 0.02$~m, $|\dot{x}| < 0.05$~m/s), while the PID drifts sinusoidally with amplitude
$\pm 0.35$~m.

\begin{figure}[H]
    \centering
    \includegraphics[width=\linewidth]{figures/position_and_speed.png}
    \caption{Horizontal position $x$ and velocity $\dot{x}$ for the three controllers.}
    \label{fig:posvel}
\end{figure}

\subsubsection{Pitch and Torque Comparison}

Figure~\ref{fig:pitchtorque} isolates the pitch and torque channels.  The MPC achieves a
smoother pitch recovery thanks to the preview capability of the MPC horizon; the SMC is faster
in the reaching phase but shows a small steady-state chattering residual of $\approx \pm 0.1^\circ$
inherent in the boundary layer.

\begin{figure}[H]
    \centering
    \includegraphics[width=\linewidth]{figures/pitch_and_torque.png}
    \caption{Pitch angle and wheel torque command.  The MPC provides the smoothest torque
    profile; the SMC shows small chattering from the boundary layer.}
    \label{fig:pitchtorque}
\end{figure}

\subsubsection{Full State Tracking and Errors}

Figure~\ref{fig:tracking_balance} provides the full state-vs-reference comparison including
CoM height.  The MPC maintains $z = 0.220$~m with sub-millimetre error; the SMC settles at
$z \approx 0.222$~m due to a small offset from the fixed posture set-points and nonlinear
kinematics; the PID oscillates by $\pm 30$~mm.

\begin{figure}[H]
    \centering
    \includegraphics[width=\linewidth]{figures/tracking_and_errors.png}
    \caption{State vs desired reference (left) and tracking errors (right): pitch, position,
    velocity, and CoM height.}
    \label{fig:tracking_balance}
\end{figure}

\subsubsection{CoM Height and Ground Contact Force}

Figure~\ref{fig:comheight} shows the CoM height and contact force.  The SMC produces a
height ripple of $\approx 3$~mm from chattering in the posture channels; the MPC tracks the
reference with $< 1$~mm error.  The PID causes repeated wheel lift-off (contact force drops
repeatedly to 0~N).

\begin{figure}[H]
    \centering
    \includegraphics[width=\linewidth]{figures/com_height_and_forces.png}
    \caption{CoM height $z$ (top) and ground contact normal force $F_n$ (bottom).  The PID
    causes wheel lift-off; MPC and SMC maintain continuous contact.}
    \label{fig:comheight}
\end{figure}

\subsection{Push Recovery Benchmark}

The PID controller is excluded from this test as it was already unstable in the free-balance
scenario.  Figure~\ref{fig:push_tracking} shows the state trajectory after the push.

\begin{figure}[H]
    \centering
    \includegraphics[width=\linewidth]{figures/tracking_and_errors.png}
    \caption{State vs reference trajectory and tracking errors after the impulsive push at
    $t = 4$~s.  MPC and SMC both recover, with different response characteristics.}
    \label{fig:push_tracking}
\end{figure}

The MPC (blue) produces a compliant response: the pitch peaks at $\approx 12^\circ$ and the robot
recovers in $\approx 3$~s via a smooth oscillatory transient dictated by the LQR cost balance.
The SMC (green) reacts more aggressively: the super-twisting drives the sliding variable back
to the chattering band in $\approx 1.5$~s, after which the residual chattering dominates.

\subsubsection{Phase Portrait and Disturbance Rejection}

Figure~\ref{fig:phase_push} shows the $(\theta,\dot\theta)$ phase portrait.  After the push
the MPC orbit follows a large, smooth arc crossing the sliding manifold; the SMC orbit
immediately seeks the manifold and slides toward the origin, consistent with second-order SMC
theory.

\begin{figure}[H]
    \centering
    \includegraphics[width=\linewidth]{figures/disturbance_and_phase_portrait.png}
    \caption{Disturbance rejection time history (left) and $(\theta,\dot\theta)$ phase portrait
    (right).  The dashed purple line is the SMC sliding manifold $s_1 = \dot\theta + c_1\theta = 0$.}
    \label{fig:phase_push}
\end{figure}

\subsubsection{Super-Twisting Sliding Variable}

Figure~\ref{fig:smc_surface} shows the sliding variable $s_1(t)$.  Before the push
($t < 4$~s) it remains within the chattering band $|s_1| < 0.1$~rad/s.  At impact it drops
sharply to $-3.5$~rad/s and recovers within $\approx 1.5$~s.  The right panel shows the state
spiralling inward toward the sliding manifold after the push.

\begin{figure}[H]
    \centering
    \includegraphics[width=\linewidth]{figures/smc_sliding_surface.png}
    \caption{Super-twisting sliding variable $s_1(t)$ (left) and SMC state trajectory in the
    $(\theta,\dot\theta)$ phase plane (right).}
    \label{fig:smc_surface}
\end{figure}

\subsection{Quantitative Comparison}

\begin{table}[H]
    \centering
    \caption{Performance summary --- free balance scenario (26~s simulation).}
    \label{tab:comparison}
    \begin{tabular}{lccc}
        \toprule
        Metric & PID & LQR+MPC & SMC \\
        \midrule
        Settling time (pitch $<0.5^\circ$) & $\infty$ (unstable) & $\approx 2.0$~s & $\approx 1.8$~s \\
        Steady-state pitch RMS & N/A & $< 0.2^\circ$ & $\approx 0.3^\circ$ \\
        Max position drift $|x|$ & $\pm 350$~mm & $< 20$~mm & $< 20$~mm \\
        CoM height error RMS & $\sim 25$~mm & $< 1$~mm & $\approx 2$~mm \\
        Max wheel torque (steady) & $\pm 6$~Nm (saturated) & $< 0.05$~Nm & $< 0.15$~Nm \\
        Contact loss events & Many & None & None \\
        Computational cost & Negligible & Medium (QP at 100~Hz) & Low (analytic at 500~Hz) \\
        Model knowledge & None & Full linear model & Nonlinear model \\
        Robustness guarantee & No & Quadratic stability & Finite-time convergence \\
        \bottomrule
    \end{tabular}
\end{table}

\subsection{Discussion}

\paragraph{PID.}  The cascade PID is conceptually simple but lacks any nonlinear model
knowledge.  Wheel torque saturation occurs because the fixed gains cannot compensate for the
coupling between the balance and posture loops far from the linearisation point.  A gain
schedule on leg length could improve the situation, but the fundamental limitation of the linear
structure remains.

\paragraph{LQR + MPC.}  The decoupled architecture provides a natural separation of concerns:
the TV-LQR handles fast balance with guaranteed quadratic stability, while the MPC previews the
CoM trajectory over a 0.5~s horizon.  The main advantage is the smooth torque profile; the main
limitation is the linear model assumption for large perturbations.

\paragraph{SMC.}  The super-twisting pitch channel provides finite-time convergence to the
sliding manifold regardless of bounded disturbances and model uncertainties.  The PFL inversion
ensures exact decoupling at the operating point.  The residual chattering visible in simulation
is a consequence of the 500~Hz discretisation and unmodelled dynamics (e.g.\ torsional resonance
of the legs); in hardware, joint compliance would attenuate it further.

The key design contribution of the SMC is precisely the choice of the \textbf{super-twisting}
algorithm for the pitch channel: it provides a continuous control law while preserving the
robustness of first-order SMC, and the integrator freeze at saturation prevents the windup that
would otherwise accumulate after a large push.

% ============================================================
\section{Contributions with Respect to Class Topics}
\label{sec:contributions}

\begin{itemize}
    \item \textbf{Wheeled robot kinematics and dynamics}: the rolling constraint is derived from
          first principles and incorporated via the actuation matrix $\bB$; the VL-WIP equivalent
          centroid follows the Lagrange--d'Alembert projection framework.
    \item \textbf{Balance and underactuated control}: three controllers span the full range from
          purely reactive (PID) to model-based predictive (MPC) to provably robust nonlinear (SMC).
    \item \textbf{Sliding mode control}: super-twisting, boundary layer, and partial feedback
          linearisation are applied to a physically meaningful multi-DOF plant, going beyond
          standard textbook examples.
    \item \textbf{Optimal control (LQR/MPC)}: TV-LQR scheduled on leg length; real-time QP at
          100~Hz in Python.
    \item \textbf{Software-in-the-loop}: full pipeline from symbolic Lagrangian (MATLAB) to
          closed-loop ROS~2/Gazebo simulation.
    \item \textbf{State estimation}: Kalman filter fusing IMU, odometry, and leg kinematics at
          500~Hz.
\end{itemize}

% ============================================================
\section{Conclusions}
\label{sec:conclusions}

This report presented the modelling and control of RoTino, a simulated VL-WIP wheeled biped
robot, covering the full pipeline from DOF analysis and Lagrangian derivation to
multi-controller benchmarking.

Key conclusions:
\begin{enumerate}
    \item The Lagrangian model captures the essential nonlinear dynamics and can be evaluated
          in real time in both MATLAB/Simulink and ROS~2/Python.
    \item The cascade PID fails to stabilise the robot in simulation due to uncompensated
          nonlinear coupling and permanent torque saturation.
    \item Both LQR+MPC and SMC achieve stable balance and disturbance rejection: the MPC gives
          the smoothest torque profile and most accurate CoM height tracking; the SMC gives
          faster reaching and stronger robustness guarantees at a lower computational cost.
    \item The super-twisting second-order SMC is identified as the best practical choice for an
          embedded real-time implementation: it is continuous, computationally inexpensive, and
          provides finite-time convergence guarantees independent of model accuracy.
\end{enumerate}

Future work could extend the model to three dimensions (full 3D yaw dynamics, inertial coupling),
add joint compliance to filter chattering, and investigate reinforcement-learning-based gain
scheduling as a function of the robot's current posture.

% ============================================================
\begin{thebibliography}{9}

\bibitem{cui2022}
Z.\ Cui, Z.\ Zhao, H.\ Zhu, J.\ Bao, and Z.\ Zhang,
``Modeling and Control of a Wheeled Biped Robot,''
\emph{Micromachines}, vol.\ 13, no.\ 5, p.\ 747, 2022.

\bibitem{levant1993}
A.\ Levant,
``Sliding order and sliding accuracy in sliding mode control,''
\emph{International Journal of Control}, vol.\ 58, no.\ 6, pp.\ 1247--1263, 1993.

\bibitem{moreno2012}
J.\ A.\ Moreno and M.\ Osorio,
``Strict Lyapunov functions for the super-twisting algorithm,''
\emph{IEEE Transactions on Automatic Control}, vol.\ 57, no.\ 4, pp.\ 1035--1040, 2012.

\bibitem{rawlings2017}
J.\ B.\ Rawlings, D.\ Q.\ Mayne, and M.\ Diehl,
\emph{Model Predictive Control: Theory, Computation, and Design}, 2nd~ed.,
Nob~Hill~Publishing, 2017.

\end{thebibliography}

\end{document}
